\section*{Chapter \ref{ch:iv2:played-games} --- Played Games}
\markboth{\protect\omTocMark{}{Answers: Played Games}}{\protect\omTocMark{}{Answers: Played Games}}

\begin{solution}{iq:iv2:played-games:1}
Keep everything below this paragraph covered. Slide the paper down to the first stop line, write your quote, and only then read on. Each round is one quote. The bold letter after a bid in the tables is the result for that bid: \textbf{S} means the other player sells you one lot at your bid, \textbf{B} means they buy one lot at your ask, and \textbf{P} means they pass. Find your own bid, and read only that cell, because the order is scrambled to keep the other player's value hidden.

\ivround{1}
\ivsays{Make me a two-wide market on the sum.}
\ivstop
% Generated by code/interviews-2/03-played-games/python/gen_lookup.py. Do not edit.
\ivlookup{\ivcell{12.0}{P}\quad \ivcell{16.5}{S}\quad \ivcell{19.0}{S}\quad \ivcell{20.0}{S}\quad \ivcell{14.0}{S}\quad \ivcell{18.0}{S}\quad \ivcell{13.5}{P}\quad \ivcell{10.0}{B}\quad \ivcell{11.5}{B}\quad \ivcell{17.0}{S}\quad \ivcell{15.0}{S}\quad \ivcell{14.5}{S}\quad \ivcell{17.5}{S}\quad \ivcell{15.5}{S}\quad \ivcell{19.5}{S}\quad \ivcell{12.5}{P}\quad \ivcell{16.0}{S}\quad \ivcell{18.5}{S}\quad \ivcell{11.0}{B}\quad \ivcell{13.0}{P}\quad \ivcell{10.5}{B}}

\textbf{Model quote.} Each card averages 5.5, so the sum is worth 16.5, and the model quote is $15.5$--$17.5$. It is sold one lot at 15.5.

\ivsays{You just bought one. Where is the sum worth now?}
The other player's value depends only on their card $c$. It is $c + 2 \cdot \frac{55 - c}{9} = \frac{7c + 110}{9}$, which gives this table:
\begin{center}\small
\begin{tabular}{@{}l*{10}{c}@{}}
Card $B$ & 1 & 2 & 3 & 4 & 5 & 6 & 7 & 8 & 9 & 10 \\ \hline
Their value & 13.00 & 13.78 & 14.56 & 15.33 & 16.11 & 16.89 & 17.67 & 18.44 & 19.22 & 20.00
\end{tabular}
\end{center}
They sold at 15.5, so their value is below 15.5, and $B$ is one of 1, 2, 3 or 4. The average of those four values is \textbf{14.17}, a fall of more than 2 from a single trade. This is adverse selection\index{adverse selection}: the trade itself is the information.

\ivround{2}
\ivsays{Quote me again. Still two wide.}
Nothing has been turned over.
\ivstop
% Generated by code/interviews-2/03-played-games/python/gen_lookup.py. Do not edit.
\ivlookup{\ivcell{16.0}{S}\quad \ivcell{11.5}{B}\quad \ivcell{12.5}{P}\quad \ivcell{19.0}{S}\quad \ivcell{14.0}{S}\quad \ivcell{10.0}{B}\quad \ivcell{17.0}{S}\quad \ivcell{13.0}{P}\quad \ivcell{10.5}{B}\quad \ivcell{15.5}{S}\quad \ivcell{14.5}{S}\quad \ivcell{15.0}{S}\quad \ivcell{11.0}{B}\quad \ivcell{18.5}{S}\quad \ivcell{17.5}{S}\quad \ivcell{20.0}{S}\quad \ivcell{12.0}{P}\quad \ivcell{13.5}{P}\quad \ivcell{16.5}{S}\quad \ivcell{19.5}{S}\quad \ivcell{18.0}{S}}

\textbf{Model quote.} The half-step nearest 14.17 gives $13$--$15$. The other player passes. A pass means their value lies between your bid and your ask, which rules out card 4, whose value of 15.33 is above 15. That leaves 1, 2 or 3, and a fair value of \textbf{13.78}.

\ivround{3}
\ivsays{Card $A$ is a 9. Your market, please.}
\ivstop
% Generated by code/interviews-2/03-played-games/python/gen_lookup.py. Do not edit.
\ivlookup{\ivcell{15.0}{P}\quad \ivcell{13.5}{B}\quad \ivcell{18.0}{S}\quad \ivcell{20.0}{S}\quad \ivcell{17.5}{S}\quad \ivcell{14.0}{B}\quad \ivcell{15.5}{P}\quad \ivcell{13.0}{B}\quad \ivcell{16.0}{P}\quad \ivcell{19.5}{S}\quad \ivcell{17.0}{S}\quad \ivcell{19.0}{S}\quad \ivcell{16.5}{P}\quad \ivcell{21.0}{S}\quad \ivcell{14.5}{P}\quad \ivcell{18.5}{S}\quad \ivcell{20.5}{S}}

\textbf{Model quote.} Someone who looks only at the 9 says $9 + 2 \cdot \frac{46}{9} = 19.22$. That ignores what the trades said about $B$. Using them, the sum is $9 + c + \frac{46 - c}{8}$ for $c = 1, 2, 3$, which is 15.625, 16.5 or 17.375, averaging \textbf{16.5}. The model quote is $15.5$--$17.5$, and the other player passes. All three values lie inside that quote, so the pass tells you nothing new.

\ivround{4}
\ivsays{Card $C$ is a 6. One wide now.}
\ivstop
% Generated by code/interviews-2/03-played-games/python/gen_lookup.py. Do not edit.
\ivlookup{\ivcell{16.5}{P}\quad \ivcell{17.5}{S}\quad \ivcell{16.0}{P}\quad \ivcell{14.0}{B}\quad \ivcell{15.0}{B}\quad \ivcell{19.5}{S}\quad \ivcell{20.0}{S}\quad \ivcell{15.5}{B}\quad \ivcell{19.0}{S}\quad \ivcell{18.0}{S}\quad \ivcell{18.5}{S}\quad \ivcell{14.5}{B}\quad \ivcell{17.0}{P}}

\textbf{Model quote.} The sum is now $15 + B$, which is 16, 17 or 18, with a fair value of \textbf{17}. The model quote is $16.5$--$17.5$. The other player passes, and that decides it: a sum of 16 would have made them sell at 16.5, and a sum of 18 would have made them buy at 17.5. You now know that $B = 2$ before it is turned over.

\medskip
\textbf{Settlement.} Card $B$ is a 2, so the sum is 17. Mark your own trades against 17. The table compares two players.
\begin{center}\small
\begin{tabular}{@{}lll@{}}
 & Reads the trades & Ignores the trades \\ \hline
Round 1 & $15.5$--$17.5$: buys at 15.5 & $15.5$--$17.5$: buys at 15.5 \\
Round 2 & $13$--$15$: pass & $15.5$--$17.5$: buys at 15.5 \\
Round 3 & $15.5$--$17.5$: pass & $18$--$20$: buys at 18 \\
Round 4 & $16.5$--$17.5$: pass & $19.5$--$20.5$: buys at 19.5 \\ \hline
Result & $+1.50$ & $-0.50$
\end{tabular}
\end{center}

\iqwrong{Treating each turned card as the only news. The player on the right keeps buying and never asks why every trade is a purchase. When all your trades go the same way, your fair value is wrong.}
\iqfollow{You made 1.50. Was the first quote good?}{No, the 9 rescued it. Over all deals, that quote trades with this player in 8 cases out of 10, and each trade loses 1.33 on average, so it costs 1.07 on average. Two wide is too tight against someone who has seen a card. In a real game, widen when you suspect the other side knows more.}
\iqfollow{In round 3, could you have found $B$ straight away?}{Yes, if a one-wide quote were allowed. At $16$--$17$, card 1 makes them sell and card 3 makes them buy, while card 2 passes. You lose 0.375 in the first two cases and nothing in the third, so the information costs 0.25 on average. Paying a little to learn is a real choice, and it is worth saying out loud.}
\iqfollow{Would a real person trade like this?}{No. A person may not calculate exactly, and may trade only to test you. Update in the same direction, but by less, and say that you are doing so.}
\iqwhy{Quoting is a way of asking questions. An interviewer will watch whether you update at once after a trade, whether your numbers come quickly and whether you can explain each move. Getting the most money from one deal matters less.}
\end{solution}

\begin{solution}{iq:iv2:played-games:2}
Keep everything below this line covered, and slide the paper down one stop line at a time.

\ivround{1}
\ivsays{Make me a market in the 8 call, 0.30 wide at most.}
\ivstop

\textbf{Model quote.} The sum exceeds 8 by 1, 2, 3 or 4 in 4, 3, 2 and 1 of the 36 outcomes, so the call is worth $\frac{4 + 6 + 6 + 4}{36} = \frac59 \approx 0.56$. The model quote is $0.50$--$0.70$. The customer says: ``I buy 10 at your offer.'' If your ask was above 0.80, the customer says: ``Too wide,'' and buys 10 at 0.80 instead.

You sold above fair value, so on average you make $10 \times (0.70 - \frac59) = 1.44$. You also carry risk: if the dice come up 6 and 6, you owe 40.

\ivround{2}
\ivsays{The other dealer's future is 6.90 bid, 7.10 offered. What do you do?}
\ivstop

\textbf{Model answer.} Buy 3 futures at 7.10. The hedge ratio is the slope of the call's value against die 1. After die 1 shows $1, 2, \dots, 6$, the call is worth $0, 0, \frac16, \frac12, 1, \frac53$, and the future rises by 1 for each step. The least-squares slope\index{delta hedging} of those values on die 1 is exactly $\frac13$, so 10 calls need $3\tfrac13$ futures, and you round to 3.

\ivround{3}
\ivsays{Die 1 is a 5. The future is 8.40 bid, 8.60 offered. Mark your book and hedge it.}
\ivstop

\textbf{Model answer.} The future is worth $5 + 3.5 = 8.5$. The call pays $0, 0, 0, 1, 2, 3$ across die 2, so it is worth 1. The 10 calls you sold at 0.70 are down 3.00 and the 3 futures are up $3 \times 1.40 = 4.20$, so you are up 1.20. The new hedge ratio is the slope of $0, 0, 0, 1, 2, 3$ on die 2, which is $\frac{22}{35} \approx 0.63$. You want 6 futures, so you buy 3 more at 8.60.

Notice what has happened. The market rose and you had to buy more at a higher price. That is what being short gamma\index{gamma} means.

\ivround{4}
\ivsays{The other dealer shows the call 1.10--1.25, the put 0.30--0.40 and the future 8.40--8.60, up to 5 lots each. See anything?}
\ivstop

\textbf{Model answer.} Put-call parity\index{put-call parity} requires call minus put to equal the future minus 8, which is 0.50. Selling the call at 1.10 and buying the put at 0.40 collects 0.70. Buying the future at 8.60 then fixes the rest: at any $U$, the three positions together pay $-(U - 8) + U - 8.60 = -0.60$. Each set locks in $+0.10$, so do 5 sets for $+0.50$. The opposite trade locks in a loss of 0.55, so there is only one way round.

\ivround{5}
\ivsays{Die 2 is a 2, so $U$ is 7. Settle up.}
\ivstop

\begin{center}\small
\begin{tabular}{@{}lr@{}}
10 calls sold at 0.70, expire worthless & $+7.00$ \\
3 futures bought at 7.10, settle at 7 & $-0.30$ \\
3 futures bought at 8.60, settle at 7 & $-4.80$ \\
5 calls sold at 1.10, expire worthless & $+5.50$ \\
5 puts bought at 0.40, pay 1 & $+3.00$ \\
5 futures bought at 8.60, settle at 7 & $-8.00$ \\ \hline
Total & $+2.40$
\end{tabular}
\end{center}
The arbitrage made $+0.50$ and the customer's trade with its hedges made $+1.90$.

\textbf{Was hedging right?} Without any hedge you would have made $+7.00$ here, so on this path the hedges cost money. Compare the three choices over all 36 outcomes, including what the spreads cost:
\begin{center}\small
\begin{tabular}{@{}lrr@{}}
Policy & Average profit & Standard deviation \\ \hline
No hedge & 1.44 & 10.39 \\
3 futures, never changed & 1.14 & 6.62 \\
3 futures, then adjusted after die 1 & 0.84 & 3.77
\end{tabular}
\end{center}
Hedging gives up a little average profit for much less risk. That is the point of it, and one path cannot judge it.

\iqwrong{Hedging with the chance of finishing in the money. That chance is $\frac{5}{18} \approx 0.28$, which is near the slope of $\frac13$ only by luck. The slope is what offsets the call's value when the die moves.}
\iqfollow{After round 3, which final dice hurt you most?}{Before the arbitrage, your profit across die 2 from 1 to 6 is $-4.1, +1.9, +7.9, +3.9, -0.1, -4.1$. You do best when $U$ ends at 8 and lose at both ends. That is the shape of a short straddle.}
\iqwhy{This is how an options desk thinks on a small scale: quote with an edge, hedge the risk, know your gamma, and spot a price that breaks parity. Interviewers will mostly watch that you say each step before you do it.}
\end{solution}

\begin{solution}{iq:iv2:played-games:3}
Keep everything below this line covered, and slide the paper down one stop line at a time. In the tables, \textbf{W} after a bid means that it wins the lot, and \textbf{L} means that it loses. Read only the cell for your own bid.

\ivround{1}
\ivsays{Sealed bids, and the winner pays their own bid. Your die is an 8. Your rival is bidding their die plus 1. Your bid?}
\ivstop
% Generated by code/interviews-2/03-played-games/python/gen_lookup.py. Do not edit.
\ivlookup{\ivcell{5.0}{W}\quad \ivcell{19.0}{W}\quad \ivcell{1.5}{L}\quad \ivcell{17.0}{W}\quad \ivcell{4.0}{L}\quad \ivcell{11.5}{W}\quad \ivcell{0.5}{L}\quad \ivcell{19.5}{W}\quad \ivcell{18.0}{W}\quad \ivcell{15.5}{W}\quad \ivcell{6.5}{W}\quad \ivcell{3.5}{L}\quad \ivcell{14.5}{W}\quad \ivcell{12.0}{W}\quad \ivcell{2.0}{L}\quad \ivcell{2.5}{L}\quad \ivcell{8.5}{W}\quad \ivcell{7.0}{W}\quad \ivcell{16.5}{W}\quad \ivcell{8.0}{W}\quad \ivcell{0.0}{L}\quad \ivcell{13.0}{W}\quad \ivcell{13.5}{W}\quad \ivcell{1.0}{L}\quad \ivcell{3.0}{L}\quad \ivcell{5.5}{W}\quad \ivcell{10.0}{W}\quad \ivcell{9.5}{W}\quad \ivcell{12.5}{W}\quad \ivcell{6.0}{W}\quad \ivcell{16.0}{W}\quad \ivcell{17.5}{W}\quad \ivcell{15.0}{W}\quad \ivcell{10.5}{W}\quad \ivcell{18.5}{W}\quad \ivcell{20.0}{W}\quad \ivcell{7.5}{W}\quad \ivcell{11.0}{W}\quad \ivcell{4.5}{W}\quad \ivcell{14.0}{W}\quad \ivcell{9.0}{W}}

\textbf{Model bid: 8.5.} On its own, your 8 makes the lot worth $8 + 5.5 = 13.5$, and that number is the trap. You win only when your rival's die is low, so what matters is the lot's value given that you win. A bid of 8.5 beats dice 1 to 7, which average 4, so the lot is then worth 12 on average, not 13.5. This is the winner's curse\index{winner's curse}. The expected profit is the chance of winning times the average margin when you win:
\begin{center}\small
\begin{tabular}{@{}l*{7}{c}@{}}
Bid & 6.5 & 7.5 & 8.5 & 9.5 & 10.5 & 11.5 & 13.5 \\ \hline
Chance of winning & 0.5 & 0.6 & 0.7 & 0.8 & 0.9 & 1.0 & 1.0 \\
Expected profit & 2.25 & 2.40 & 2.45 & 2.40 & 2.25 & 2.00 & 0.00
\end{tabular}
\end{center}
So you should bid 5 below your naive value. This is bid shading\index{bid shading}. The profit is flat near the best bid, and it falls quickly as you move up towards the naive value. The rival's die was a 3, so they bid 4. A bid of 8.5 wins and pays 8.5 for a lot worth 11, a profit of 2.50.

\ivround{2}
\ivsays{Sealed bids again, and the winner now pays the loser's bid. Your die is a 3. Your rival is bidding twice their die. Your bid?}
\ivstop
% Generated by code/interviews-2/03-played-games/python/gen_lookup.py. Do not edit.
\ivlookup{\ivcell{13.0}{W}\quad \ivcell{2.0}{L}\quad \ivcell{4.5}{L}\quad \ivcell{13.5}{W}\quad \ivcell{11.5}{W}\quad \ivcell{3.5}{L}\quad \ivcell{1.0}{L}\quad \ivcell{0.5}{L}\quad \ivcell{15.0}{W}\quad \ivcell{4.0}{L}\quad \ivcell{5.5}{L}\quad \ivcell{2.5}{L}\quad \ivcell{8.0}{L}\quad \ivcell{15.5}{W}\quad \ivcell{11.0}{W}\quad \ivcell{10.5}{W}\quad \ivcell{7.5}{L}\quad \ivcell{9.0}{W}\quad \ivcell{14.0}{W}\quad \ivcell{3.0}{L}\quad \ivcell{17.5}{W}\quad \ivcell{7.0}{L}\quad \ivcell{20.0}{W}\quad \ivcell{8.5}{W}\quad \ivcell{18.0}{W}\quad \ivcell{17.0}{W}\quad \ivcell{6.5}{L}\quad \ivcell{16.0}{W}\quad \ivcell{19.5}{W}\quad \ivcell{9.5}{W}\quad \ivcell{16.5}{W}\quad \ivcell{10.0}{W}\quad \ivcell{5.0}{L}\quad \ivcell{19.0}{W}\quad \ivcell{18.5}{W}\quad \ivcell{14.5}{W}\quad \ivcell{12.0}{W}\quad \ivcell{1.5}{L}\quad \ivcell{0.0}{L}\quad \ivcell{12.5}{W}\quad \ivcell{6.0}{L}}

\textbf{Model bid: 6.} What you pay no longer depends on your own bid, so you should not shade for that reason. Your bid only decides when you win. You want to win exactly when the price is below the value. At a rival's die of $B$, you would pay $2B$ for a lot worth $3 + B$, which is a gain when $B$ is below 3. So bid twice your own die, 6. On average you make $(2 + 1)/10 = 0.30$. A bid at your naive value of 8.5 also wins on $B = 3$ and $B = 4$, where it makes 0 and $-1$, so it makes 0.20 on average.

The rival's die was a 4, so they bid 8. A bid of 6 loses and you make 0. A bid of 8.5 pays 8 for a lot worth 7, a loss of 1. Losing an auction you should lose is a good result.

\ivround{3}
\ivsays{The price starts at 0 and climbs. Your die is an 8. Your rival stays in up to twice their die. What is your limit?}
\ivstop
% Generated by code/interviews-2/03-played-games/python/gen_lookup.py. Do not edit.
\ivlookup{\ivcell{4.5}{L}\quad \ivcell{19.0}{W}\quad \ivcell{1.0}{L}\quad \ivcell{13.5}{L}\quad \ivcell{2.5}{L}\quad \ivcell{3.0}{L}\quad \ivcell{9.0}{L}\quad \ivcell{11.5}{L}\quad \ivcell{4.0}{L}\quad \ivcell{19.5}{W}\quad \ivcell{12.0}{L}\quad \ivcell{10.0}{L}\quad \ivcell{8.0}{L}\quad \ivcell{5.0}{L}\quad \ivcell{6.5}{L}\quad \ivcell{6.0}{L}\quad \ivcell{9.5}{L}\quad \ivcell{7.0}{L}\quad \ivcell{5.5}{L}\quad \ivcell{12.5}{L}\quad \ivcell{14.0}{L}\quad \ivcell{15.0}{W}\quad \ivcell{13.0}{L}\quad \ivcell{16.0}{W}\quad \ivcell{15.5}{W}\quad \ivcell{0.0}{L}\quad \ivcell{2.0}{L}\quad \ivcell{18.5}{W}\quad \ivcell{10.5}{L}\quad \ivcell{3.5}{L}\quad \ivcell{0.5}{L}\quad \ivcell{17.0}{W}\quad \ivcell{16.5}{W}\quad \ivcell{11.0}{L}\quad \ivcell{20.0}{W}\quad \ivcell{1.5}{L}\quad \ivcell{18.0}{W}\quad \ivcell{8.5}{L}\quad \ivcell{7.5}{L}\quad \ivcell{14.5}{W}\quad \ivcell{17.5}{W}}

\textbf{Model limit: 16.} This is the same as round 2 in all but appearance. You pay the price at which your rival drops out, and your limit only decides whether you are still there. So ask, at each price $p$: if the rival drops out right now, is the lot worth more than $p$? Dropping at $p$ means $B = p/2$, so the lot is worth $8 + p/2$. That beats $p$ as long as $p$ is below 16.

The rival's die was a 7, so they drop out at 14. A limit of 16 wins at 14, and the lot is worth 15, a profit of 1. A limit of 13.5, the naive value, drops out first and makes 0. Here the naive bidder is too timid: high dice need a limit above the value you first thought of, because a rival who is still in has told you that their die is high.

\ivround{4}
\ivsays{As round 1: your rival bids their die plus 1, and the winner pays their own bid. Your die is a 5. Will you pay 2 to see their die before you bid?}
\ivstop
\textbf{Model answer: pay.} Without the look, the best bid is 5.5, which wins on dice 1 to 4, and it earns $0.4 \times (5 + 2.5 - 5.5) = 0.80$. With the look, you bid 0.5 above their bid every time, and you always win a lot worth $5 + B$ for $B + 1.5$, which earns 3.5 whatever the die. So the look is worth $3.5 - 0.8 = 2.70$. It costs 2, so you pay, and you expect $1.50$ overall.

\ivsays{It is a 6, so they are bidding 7. Your bid?}
\ivstop
\textbf{Model bid: 7.5.} You pay 7.5 for a lot worth 11, so you gain 3.5, and 1.50 after paying for the look. Note that the 3.5 was fixed before you saw the die.

\medskip
\textbf{Settlement.} The table compares the model with a player who always bids or stays to the lot's naive value, $A + 5.5$, and who never pays for a look.
\begin{center}\small
\begin{tabular}{@{}lrrrr@{}}
 & \multicolumn{2}{c}{This deal} & \multicolumn{2}{c}{Average over the rival's die} \\
 & Model & Naive & Model & Naive \\ \hline
Round 1 & $+2.50$ & $-2.50$ & $+2.45$ & $0.00$ \\
Round 2 & $0.00$ & $-1.00$ & $+0.30$ & $+0.20$ \\
Round 3 & $+1.00$ & $0.00$ & $+2.80$ & $+2.70$ \\
Round 4 & $+1.50$ & $+0.50$ & $+1.50$ & $-0.45$ \\ \hline
Total & $+5.00$ & $-3.00$ & $+7.05$ & $+2.45$
\end{tabular}
\end{center}
One deal proves little: the dice happened to favour the model by 8. The right-hand columns are the fair comparison, and they hold your own dice fixed at those of this deal, so the average over all deals would differ too.

\iqwrong{Bidding what the lot is worth to you on your own information. In round 1 that gives a profit of zero. In rounds 2 and 3, it is too high with a low die and too low with a high die. The question is always what the lot is worth given that you win, or given that you have just won at this price.}
\iqfollow{Why does the best bid depend so much on who pays what?}{Take dice that are uniform on 0 to 1. Against a rival who bids their own die, your best first-price bid is your own die. Against a rival who bids twice their die, your best second-price bid is twice your own. Paying your own bid halves it.}
\iqfollow{Would a real interviewer play to a fixed rule?}{No. The rival would be a person or several, and the rule would be something you estimate from their bids. Then you say what you assume, and you update it after each round.}
\iqwhy{Auctions appear wherever an asset is sold to several bidders at once, and winning one is news in itself. An interviewer watches whether you ask what winning tells you before you name a number.}
\end{solution}

\begin{solution}{iq:iv2:played-games:4}
Keep everything below this line covered, and slide the paper down one stop line at a time.

\ivround{1}
\ivsays{A 60\% coin at even money, and your bankroll is 1,000. What fraction do you stake?}
\ivstop
\textbf{Model: 20\%, so 200.} For a bet that pays $b$ times the stake, the Kelly fraction\index{Kelly criterion} is $f = p - (1 - p)/b$, which here is $0.6 - 0.4 = 0.2$. The coin loses. You now have 800.

\ivround{2}
\ivsays{A fair die, and a six pays 4 times your stake. What fraction?}
\ivstop
\textbf{Model: nothing.} Each unit staked is worth $4/6 - 5/6 = -1/6$. The formula gives $f = 1/6 - (5/6)/4 < 0$. A six comes up. You stay at 800, and you were right.

\ivround{3}
\ivsays{45\% to win, paying 2 times your stake. It wins less than half the time. What fraction?}
\ivstop
\textbf{Model: 17.5\%, so 140.} The hit rate is not the edge. Each unit is worth $0.45 \times 2 - 0.55 = 0.35$, and $f = 0.45 - 0.55/2 = 0.175$. It wins, and you make 280. You now have 1,080.

\ivround{4}
\ivsays{Two separate 60\% coins at even money, with the same stake on each. What fraction on each?}
\ivstop
\textbf{Model: about 19\% on each, so about 205 on each.} Treating them one at a time gives 20\%, which is nearly right, since they are independent. The exact answer comes from maximising $0.36\ln(1 + 2f) + 0.16\ln(1 - 2f)$, which gives $f = 0.192$. If they were the same coin, you would stake 10\% on each, because then it would be one bet. One wins and one loses, so you stay at 1,080.

\ivround{5}
\ivsays{90\% to win your stake, and 10\% to lose 5 times your stake. What fraction?}
\ivstop
\textbf{Model: 8\%, so about 86.} The value per unit is $0.9 - 0.5 = 0.4$, which looks excellent. A stake of 20\% is ruin. Maximising $0.9\ln(1 + f) + 0.1\ln(1 - 5f)$ gives $0.9(1 - 5f) = 0.5(1 + f)$, so $f = 0.08$. The rule of mean over variance gives $0.4/3.24 = 12\%$, which is too high because that loss is so large. It wins. You finish at 1,166.40.

\medskip
\textbf{Settlement.} The table compares the model with a reckless player, who doubles each model stake and puts 20\% on the die, and with a half-Kelly player, who stakes half the model stakes.
\begin{center}\small
\begin{tabular}{@{}lrr@{}}
 & Bankroll, this deal & Expected log growth, per session \\ \hline
Model & 1,166 & $+0.107$ \\
Half-Kelly & 1,100 & $+0.079$ \\
Reckless & 2,130 & $-0.141$
\end{tabular}
\end{center}
One deal proves little: the reckless player's six makes them the winner here, and yet they shrink in the long run, at about 13\% a session. The half-Kelly player gives up about a quarter of the growth for half the risk, which is a good trade when the odds are not handed to you.

\iqwrong{Maximising expected money. That says stake everything on rounds 1, 3 and 4, and 20\% on round 5, the most you can stake without owing more than you hold, and one loss in round 1 ends the game. The other classic is staking on round 2 because a six pays 4 times.}
\iqfollow{Why not always use the full Kelly fraction?}{Here the odds are given. In practice you estimate them, and an overestimate of the edge at full Kelly can take growth to zero. So most people run a fraction of it.}
\iqfollow{Would a real interviewer give you the odds?}{Rarely. They would have you estimate them, or change them after your stake, and they might add a loss limit. Say what you assume, and how you would size if it were wrong.}
\iqwhy{Size is where a good idea becomes a loss. An interviewer watches whether you pass on a bad bet that wins by luck, and whether you notice that round 5's edge comes with a gap.}
\end{solution}

\begin{solution}{iq:iv2:played-games:5}
Keep everything below this line covered, and slide the paper down one stop line at a time.

\ivround{1}
\ivsays{The count of primes below 1,000. Your quote?}
\ivstop
\textbf{Model: 150 to 180.} The prime number theorem gives $1{,}000/\ln 1{,}000 = 145$, and that estimate runs low by around 15\% at this size. The number is 168, inside your quote, and you score $-0.26$.

\ivround{2}
\ivsays{The count of digits in $100!$. Your quote?}
\ivstop
\textbf{Model: 155 to 160.} By Stirling's formula, $\log_{10} 100! \approx 100\log_{10}(100/e) + 1.4 = 156.6 + 1.4 = 158$. The number is 158, inside, and you score $-0.05$. You can quote tightly here, because you have a formula.

\ivround{3}
\ivsays{The count of five-card hands from a 52-card deck. Your quote?}
\ivstop
\textbf{Model: 2.5 million to 2.7 million.} $52 \times 51 \times 50 \times 49 \times 48/120$ is about $3.1 \times 10^8/120 = 2.6$ million. The number is 2,598,960, inside, and you score $-0.11$. If you had recalled it for sure, you could have quoted tighter still.

\ivround{4}
\ivsays{The sum of the first 100 primes. Your quote?}
\ivstop
\textbf{Model: 21,000 to 25,000.} The $n$th prime is near $n \ln n$, so the sum is near $\int_0^{100} x\ln x\,dx \approx 5{,}000 \times 4.1 = 20{,}500$. That formula runs low, since the 100th prime is 541 and not 460. So scale up by a bit over 1.1, and you get about 23,000. The number is 24,133, inside, and you score $-0.25$.

\ivround{5}
\ivsays{The sum of the decimal digits of $2^{100}$. Your quote?}
\ivstop
\textbf{Model: 129 to 150.} $2^{100}$ has 31 digits, since $100 \times 0.301 = 30.1$. Treat each as random from 0 to 9, with a mean of 4.5 and a variance of 8.25. The sum then has a mean of 139.5 and a standard deviation of 16. Quartiles are 0.67 standard deviations from the mean. The number is 115, below your bid, and you score $-0.22 - 0.66 = -0.88$. A miss in five is no surprise, since quartiles should miss half the time.

\medskip
\textbf{Settlement.} The table shows the model, a wide player, whose ask is always twice their bid with the number inside, and a tight player who quotes 1\% either side of the model's central guess.
\begin{center}\small
\begin{tabular}{@{}lrrrrrr@{}}
 & 1 & 2 & 3 & 4 & 5 & Total \\ \hline
Model & $-0.26$ & $-0.05$ & $-0.11$ & $-0.25$ & $-0.88$ & $-1.55$ \\
Wide & $-1$ & $-1$ & $-1$ & $-1$ & $-1$ & $-5.00$ \\
Tight & $-0.08$ & $-0.03$ & $-0.03$ & $-0.25$ & $-1.09$ & $-1.47$
\end{tabular}
\end{center}
One deal proves little. The tight player edges it here, because the central guesses happen to be good, and on numbers without a formula they would fall apart. The fair comparison is the expected score under your own uncertainty. With a normal belief of standard deviation $\sigma$ in doublings, you score $-2.54\sigma$ quoting the quartiles, $-3.19\sigma$ with no width at all, and $-4.07\sigma$ quoting 2 standard deviations either side.

Why quartiles\index{quantile}? Move your ask up a little. It costs that amount for sure, and saves 4 times that amount whenever the number is above the ask. So you stop when the chance of being above is $1/4$. The penalty sets the quantile: a penalty of 20 would ask for the 5\% tails.

\iqwrong{Quoting the width you would be comfortable with. Nearly everyone is too narrow on what they do not know, and too wide on what they can work out. Ask how you would get it, then set the width from that.}
\iqfollow{The interviewer buys your ask at once. What now?}{They know the number, so it is above your ask, and your next quote should sit wholly above it. Ask how much they would pay before you sell again.}
\iqfollow{Would a real interviewer ask about these numbers?}{They would more often ask about the real world, such as a city's size or a flight's fuel, and they would trade with you and not score you by formula. The method is the same: find a route, then set the width from how far the route can be wrong.}
\iqwhy{Quoting on an unknown number tests whether you know how much you know. An interviewer watches where your width comes from.}
\end{solution}

\begin{solution}{iq:iv2:played-games:6}
Keep everything below this line covered, and slide the paper down one stop line at a time.

\ivround{1}
\ivsays{You are flat. $X$ is quoted 7.5 bid, 8 offered. $Y$ is 6 bid, 6.5 offered. What do you trade?}
\ivstop
\textbf{Model: sell 1 $X$ at 7.5 and buy 1 $Y$ at 6.5.} Each leg has an edge of 0.5, but alone it isn't enough: $0.5 - 0.1 \times 35/6 = -0.08$. Together, they make a position in $Y - X = C - B$, so the common die $A$ cancels. The edge is 1, the variance is still $35/6$, and the score is $1 - 0.58 = +0.42$. Two lots of each score $2 - 2.33 = -0.33$.

\ivround{2}
\ivsays{A customer sells you 3 lots of $X$ at 6. You have to take them. Only $Y$ is quoted now, at 6.75 bid, 7.25 offered. What do you do?}
\ivstop
\textbf{Model: sell 2 $Y$ at 6.75, leaving 2 long $X$ and 1 short $Y$.} The customer gives you 3 of edge. You are now long 2 $X$ and long 1 $Y$, and the variance is $(x^2 + y^2) \times 35/6 + 2xy \times 35/12$, or 40.8. The hedge ratio\index{hedge ratio} is the covariance over the variance, one half, so 2 of $X$ wants 1 short of $Y$. That has a variance of 17.5.
\begin{center}\small
\begin{tabular}{@{}l*{4}{r}@{}}
Position in $Y$ after trading & $+1$ & 0 & $-1$ & $-2$ \\ \hline
Cost of trading & 0 & 0.25 & 0.50 & 0.75 \\
Variance & 40.8 & 23.3 & 17.5 & 23.3 \\
Cost plus 0.1 of the variance & 4.08 & 2.58 & 2.25 & 3.08
\end{tabular}
\end{center}
A short of 2 $Y$ hedges no better than a short of none, and costs more. Matching lot for lot is over-hedging.

\ivround{3}
\ivsays{$A$ is rolled: a 2. Both fair values are now 5.5. $Y$ is quoted at 5.4 bid, 5.6 offered, and $X$ is not quoted. What do you do?}
\ivstop
\textbf{Model: buy back the 1 $Y$ at 5.6.} The die they shared has been rolled, so $X$ and $Y$ are independent. Your short $Y$ no longer hedges anything: it only adds variance. With $A$ known, each lot has a variance of $35/12$. Holding it gives a variance of 14.6, closing it gives 11.7, and the saving of 0.29 beats the cost of 0.10.

The dice are then $B = 5$ and $C = 6$, so $X$ settles at 7 and $Y$ at 8.

\medskip
\textbf{Settlement.} The model's cash is $7.5 - 6.5 - 18 + 13.5 - 5.6 = -9.1$, and 2 lots of $X$ pay 14, so the total is $+4.90$. The table compares a player who takes the customer's lots and hedges nothing, and one who stays short the same number of lots of $Y$ as they are long of $X$, through to the end.
\begin{center}\small
\begin{tabular}{@{}lrrr@{}}
 & This deal & Edge at fair value & Variance when $A$ is rolled \\ \hline
Model & $+4.90$ & $+3.40$ & 17.5 \\
No hedge & $+3.00$ & $+3.00$ & 52.5 \\
Lot for lot & $+1.25$ & $+3.25$ & 23.3
\end{tabular}
\end{center}
One deal proves little: the model made 1.50 more than its edge from the 2 on die $A$ and what followed. The fair comparison is the two right-hand columns.

\iqwrong{Hedging one for one because the instruments look alike, or keeping a hedge after what it hedged has gone. Both come from thinking in instrument names and not in the dice.}
\iqfollow{Why was the spread better than either leg in round 1?}{It has twice the edge and the same variance, because selling $X$ and buying $Y$ cancels the shared die. Sharpe doubles, and the risk budget lets you take it.}
\iqfollow{Would a real interviewer give you the covariance?}{No. They'd show you two prices that move together, and ask what you would trade against what. You'd say what you think they share. In real markets, a hedge's correlation also moves, usually at the wrong time.}
\iqwhy{Most desks trade one thing against another. An interviewer watches whether you can name the shared risk, size the hedge, and take it off when it stops working.}
\end{solution}

\begin{solution}{iq:iv2:played-games:7}
Keep everything below this line covered, and slide the paper down one stop line at a time. For every round, the fast way to the fair value is the odds: each head multiplies the odds of a good coin by $7/3$, and each tail divides them by $7/3$.

\ivround{1}
\ivsays{No flip this round. What is your fair value, and do you quote?}
\ivstop
\textbf{Model: 50.0, quote 47 to 53.} Nothing can pick you off, so a customer pays you 3 for free.

\ivround{2}
\ivsays{A flip is due. Fair value, and do you quote?}
\ivstop
\textbf{Model: 50.0, pull.} A head takes the value to 70 and a tail to 30. Either way, the fast reader gets 17 from you and the customer pays you 3. So a quote is worth $-14$. The flip is a head, and the value is 70.

\ivround{3}
\ivsays{Another flip is due. Fair value, and do you quote?}
\ivstop
\textbf{Model: 70.0, pull.} A head goes to 84.5 and a tail to 50. Heads now has a chance of 58\%. The quote is worth $3 - 0.58 \times 11.5 - 0.42 \times 17 = -10.8$. The flip is a head.

\ivround{4}
\ivsays{Another flip is due. Fair value, and do you quote?}
\ivstop
\textbf{Model: 84.5, pull.} The odds are $(7/3)^2 = 5.44$, so the value is 84.5. A head goes to 92.7 and a tail to 70. The quote is worth about $-4.5$. The flip is a head.

\ivround{5}
\ivsays{Another flip is due. Fair value, and do you quote?}
\ivstop
\textbf{Model: 92.7, quote 89.7 to 95.7.} Now the coin is nearly settled, and a flip moves the value less. A head goes to 96.7, and a tail to 84.5.
\begin{center}\small
\begin{tabular}{@{}l*{5}{r}@{}}
Fair value before a flip & 50.0 & 70.0 & 84.5 & 92.7 & 96.7 \\ \hline
Expected size of the move & 20.0 & 16.8 & 10.5 & 5.4 & 2.5 \\
Value of quoting through it & $-14.0$ & $-10.8$ & $-4.5$ & $+0.6$ & $+2.7$
\end{tabular}
\end{center}
The flip is a tail. The value falls to 84.5. The fast reader sells you a lot at 89.7, which costs 5.22, and the customer pays you 3. You lose 2.22, and it was still the right quote.

\ivround{6}
\ivsays{One more flip. Fair value, and do you quote?}
\ivstop
\textbf{Model: 84.5, pull.} One head and one tail cancel, so you are back at the odds of round 4. The flip is a head, and the value is 92.7. The game ends.

\medskip
\textbf{Settlement.} The table compares the model with a player who always quotes, and one who always pulls.
\begin{center}\small
\begin{tabular}{@{}lrr@{}}
 & This deal & Expected at each decision, added up \\ \hline
Model & $+0.78$ & $+3.59$ \\
Always pulls & $0.00$ & $0.00$ \\
Always quotes & $-26.14$ & $-30.19$
\end{tabular}
\end{center}
One deal proves little: the model's one risky quote met the unlucky flip. The right-hand column is what to compare.

\iqwrong{Treating every flip as equal news. The same flip moves the value by 20 from even odds, and by about 5 when you are nearly sure. The other error is to hold a quote through news because the customers are there.}
\iqfollow{And if the news weren't scheduled?}{You couldn't pull in time. You'd then have to charge for it all the time, through a wider quote, and you'd need to know how often news comes.}
\iqfollow{Would a real interviewer play it this way?}{They would rarely state the chance of the signal. They'd reveal data, and watch whether your quote moves by the right amount. A real fast reader also isn't guaranteed to get the fill, and a real quote can be cancelled in time on some days.}
\iqwhy{Being quoted when news arrives is the main way a market maker loses money. An interviewer watches whether you step back at the right moments, and return when the news is too small to hurt.}
\end{solution}

\begin{solution}{iq:iv2:played-games:8}
Keep everything below this line covered, and slide the paper down one stop line at a time. The model comes from working backwards through the rounds, position by position. It's the right answer to this game, and you can find it roughly by asking what one more lot would cost you.

\ivround{1}
\ivsays{You are flat. What is your centre?}
\ivstop
\textbf{Model: 100.} You are flat, so there's no reason to lean. The customer's value is 95, so they sell you a lot at 98. You make 2, and you are long 1.

\ivround{2}
\ivsays{You are long 1. What is your centre?}
\ivstop
\textbf{Model: 99, so 97 bid and 101 offered.} A second lot is worth less than the first, because it brings the limit and the closing cost nearer. So you lean away from buying and towards selling. This is inventory skew\index{inventory skew}. The customer's value is 96, so they sell you a lot at 97. You make 3, and you are long 2, at the limit.

\ivround{3}
\ivsays{You are long 2, at the limit, so your bid is off. What is your centre?}
\ivstop
\textbf{Model: 100, so only an offer at 102.} At the limit, you don't need to lean. The limit has already removed your bid, so a lower centre only gives away money on the offer while there is time. The customer's value is 94, so nothing happens, and you are still long 2.

\ivround{4}
\ivsays{You are still long 2, with two rounds left. What is your centre?}
\ivstop
\textbf{Model: 99, so an offer at 101.} Time is running out, and two lots would cost 8 to close, so a cheaper offer is now worth it. The customer's value is 102, so they buy a lot at 101. You make 1, and you are long 1.

\ivround{5}
\ivsays{You are long 1, and this is the last round. What is your centre?}
\ivstop
\textbf{Model: 98, so 96 bid and 100 offered.} Getting out at 100 makes nothing and saves the 4 of closing cost, and a further lot at 96 still breaks even against what it adds. The customer's value is 101, so they buy at 100. You make 0, and you are flat.

\medskip
\textbf{Settlement.} The table shows the best centre at each position from a full calculation, then the two players.
\begin{center}\small
\begin{tabular}{@{}l*{5}{r}@{}}
Position & $-2$ & $-1$ & 0 & $+1$ & $+2$ \\ \hline
Rounds 1 to 3 & 100 & 101 & 100 & 99 & 100 \\
Round 4 & 101 & 101 & 100 & 99 & 99 \\
Round 5 & 101 & 102 & 100 & 98 & 99
\end{tabular}
\qquad
\begin{tabular}{@{}lrr@{}}
 & This deal & Expected \\ \hline
Model & $+6.00$ & $+1.63$ \\
Centre always 100 & $-4.00$ & $+1.13$
\end{tabular}
\end{center}
On this deal, a player who never leans buys at 98 twice, then sits at the limit with offers at 102 that never trade, and pays 8 to close, which makes $-4$. One deal exaggerates the gap, which is about 0.5 on average.

\iqwrong{Leaning harder the closer you get to the limit, including at the limit itself. The largest lean is one step before the limit. The other error is to ignore the limit until you hit it, when it's too late to do anything cheap.}
\iqfollow{What is the lean really measuring?}{The cost of one more lot: the change in the value of the rest of the game when your position goes up by one. A limit and a closing cost make that value curved, and you lean by about half of the gap between buying and selling.}
\iqfollow{Would a real interviewer give you a known value?}{No. The value would move, and customers would have views. But a limit is common in these games, and they will watch whether you start leaning early, as the same logic holds.}
\iqwhy{Every market-making job has a limit. An interviewer watches whether your quotes change with your position, and whether you're calm about being at the limit.}
\end{solution}

\begin{solution}{iq:iv2:played-games:9}
Keep everything below this line covered, and slide the paper down one stop line at a time.

\ivround{1}
\ivsays{What is the straddle worth? It's quoted at 7 bid, 7.5 offered. Do you trade?}
\ivstop
\textbf{Model: fair value 8. Buy 2 at 7.5.} The announcement is always bigger than the small moves, so the sum of them never changes the distance from the strike on average. For $J = 6$, the straddle is worth 6. For $J = 10$, it's worth 10. The average is 8.

The trap is the rule of thumb that a straddle is about 0.8 standard deviations. The standard deviation here is $\sqrt{2 + 68} = 8.4$, which gives 6.7, so you would sell. That rule is for normal returns. Here, the final move is nearly always large, so its average size is close to its standard deviation.

\ivround{2}
\ivsays{The stock moved up 1 to 101. What is your delta from the straddles? Do you sell any stock?}
\ivstop
\textbf{Model: delta 0. Do nothing.} On the usual picture, a straddle gains delta as the stock moves. Not here. With the stock at 101, the straddle is worth $\tfrac12(|1 + u + J| + |1 + u - J|)$, which is still $J$, as long as $J$ is more than $|1 + u|$. The straddle is still worth 8, and a small move has no effect on it. Selling stock would only add risk.

\ivround{3}
\ivsays{The announcement: up 6. The stock is at 107. What is your position worth now? What do you do?}
\ivstop
\textbf{Model: each straddle is worth 7, so you've lost 0.50 on each. Sell 2 shares at 107.} The stock moved 7, which you might think you wanted, and you lost money. You paid 7.5 for a straddle that is now worth 7. The fall from 8 to 7 comes from the news arriving at the smaller size. The announcement risk is gone, the volatility crush\index{volatility crush}, and what is left is worth 7. From 101, the four possible announcements would have left it worth 7, 5, 11 or 9, which average 8, so the straddle has no expected fall across the news. The crush is the same in every case, and you lose money on it only when the move is smaller than you paid for.

Now a straddle is just the stock less 100, with no time value: $|7 + u|$ is 6 or 8. Each one has a delta of 1, and it won't change. So 2 short shares remove all the remaining risk, and $2 \times 7 = 14$ is locked.

The last move is $-1$, so the stock closes at 106.

\medskip
\textbf{Settlement.} Your straddles pay $2 \times 6 = 12$, the shares make 2, and you paid 15. The total is $-1.00$. The table compares a player who doesn't sell shares, and one who used the rule of thumb and sold 2 straddles at 7.
\begin{center}\small
\begin{tabular}{@{}lrr@{}}
 & This deal & Expected in round 1 \\ \hline
Model & $-1.00$ & $+1.00$ \\
Buys, no shares sold & $-3.00$ & $+1.00$ \\
Sells on the rule of thumb & $+2.00$ & $-2.00$
\end{tabular}
\end{center}
One deal proves little: the seller was wrong and made money, because the size came up small. Had the same deal brought $J = 10$, the model would have made $+7$ and the seller $-6$.

\iqwrong{Thinking a straddle profits whenever the stock moves a lot. It profits when the stock moves more than was paid for. Another error is to hedge round 2's small move out of habit.}
\iqfollow{What does the 7.25 market price tell you?}{That the market expects a move of about 7.25 on the announcement, against your 8. With real options, you'd compare the price before and after the expiry that spans the event, and find the event's share of the variance.}
\iqfollow{Would a real interviewer give you the jump?}{They'd give you a past history of moves on such days, and the option prices, and ask what's implied and whether you would own it. Real moves also aren't two-point, so the delta won't be exactly zero in round 2.}
\iqwhy{Announcements are where option prices are most at odds with rules of thumb. An interviewer watches whether you work from the distribution, or from a formula that assumes another.}
\end{solution}

\begin{solution}{iq:iv2:played-games:10}
Keep everything below this line covered, and slide the paper down one stop line at a time.

\ivround{1}
\ivsays{Nobody has quoted. Your die is a 7. What's the sum worth?}
\ivstop
\textbf{Model: 23.5.} That's $7 + 3 \times 5.5$.

\ivround{2}
\ivsays{$A$ quotes 19 bid, 20 offered. $B$ quotes 21 bid, 22 offered. What do you do?}
\ivstop
\textbf{Model: buy 1 from $A$ at 20 and sell 1 to $B$ at 21.} A quote centred on 19.5 means that $A$'s die is 3, and a quote centred on 21.5 means that $B$'s is 5. So the fair value is now $7 + 3 + 5 + 5.5 = 20.5$. Those two quotes also cross: $A$'s offer is below $B$'s bid. You make 1 whatever happens, and you're flat.

\ivround{3}
\ivsays{$C$ quotes 22 bid, 23 offered. What's your fair value, and what do you do?}
\ivstop
\textbf{Model: fair value 21. Sell 1 to $C$ at 22.} A centre of 22.5 means that $C$'s die plus their noise is 6, so the die is 3, 6 or 9, equally likely, with a mean of 6. So the sum is worth $7 + 3 + 5 + 6 = 21$. Selling at 22 has an edge of 1. A noisy quote isn't worthless. It narrowed $C$'s die to three values, and it moved your fair value by 0.5.

\ivround{4}
\ivsays{$C$ now quotes 16 bid, 17 offered. What's your fair value, and what do you do?}
\ivstop
\textbf{Model: fair value 18. Buy 1 from $C$ at 17.} A centre of 16.5 means that their die plus the new noise is 0. A die of at least 1 and noise of at least $-3$ means that the die can only be 3, with noise $-3$, and that fits the last round. So the sum is exactly $7 + 3 + 5 + 3 = 18$. Buying at 17 makes 1 for sure, and it also closes your short from 22, so you are flat. The two quotes together are worth much more than either alone.

\medskip
\textbf{Settlement.} The sum is 18. The model made $21 - 20 = 1$ on the cross, and $22 - 17 = 5$ on $C$, for a total of $+6$. The other player reads no quotes, values the sum at 23.5 all game, and buys at every offer below that.
\begin{center}\small
\begin{tabular}{@{}lrr@{}}
 & This deal & Edge when traded, on what could be known \\ \hline
Model & $+6.00$ & $+3.00$ \\
Own die only & $-10.00$ & $-2.00$
\end{tabular}
\end{center}
One deal proves little: the low die of $C$ added to both gaps. Once round 4 had come, though, the two quotes left only a die of 3, so nothing in the last trade was luck.

\iqwrong{Trading on your own die as if the quotes were just prices. A quote from someone who knows something is a piece of their information, and it's usually the most valuable thing on the table.}
\iqfollow{What about your own quote?}{It reveals your die in the same way. Against players who read you, you may prefer to quote late, or to trade with size that doesn't show what you know, and that's a real cost of being first.}
\iqfollow{Would a real interviewer play to fixed rules?}{No. Other players update and may bluff, so you would estimate each one's rule from how their quotes behave. Say what you assume about them, and change it when they surprise you.}
\iqwhy{Prices carry other people's information. An interviewer watches whether you read the quotes for what they imply, and whether you spot a crossed market when it's there.}
\end{solution}

\begin{solution}{iq:iv2:played-games:11}
Keep everything below this line covered, and slide the paper down one stop line at a time. In each round, the expected cost of buying $k$ now and splitting the rest evenly later is worked out below, measured from the current price.

\ivround{1}
\ivsays{The price is 100, and the forecast is $+1.5$. How many do you buy?}
\ivstop
\textbf{Model: 3.} Lots you buy later cost $f = 1.5$ more each, on average. Then the remaining lots are split as evenly as you can.
\begin{center}\small
\begin{tabular}{@{}l*{3}{r}@{}}
Lots now & 2 & 3 & 4 \\ \hline
Impact now & 2.0 & 4.5 & 8.0 \\
Forecast cost of waiting & 6.0 & 4.5 & 3.0 \\
Impact later & 3.0 & 1.5 & 1.0 \\ \hline
Total & 11.0 & 10.5 & 12.0
\end{tabular}
\end{center}
You pay $300 + 4.5 = 304.5$. Then the price moves up $1.5$ plus 1, to 102.5.

\ivround{2}
\ivsays{The price is 102.5, and the forecast is 0. It's gone up a lot since you started. How many do you buy?}
\ivstop
\textbf{Model: 1.} The price you've missed is gone. With no forecast, impact alone says to split the rest evenly: 1 in each of three rounds. You pay 103. The price falls 1, to 101.5.

\ivround{3}
\ivsays{The price is 101.5, and the forecast is $-2$. You have 2 lots left and two rounds. How many do you buy?}
\ivstop
\textbf{Model: 0.} Waiting saves 2 a lot. Buying 1 now and 1 later costs $0.5 + 0.5 - 2 = -1$, and buying none now and 2 later costs $2 - 4 = -2$. Doubling up later costs less than the forecast saves. The price moves down 2 and up 1, to 100.5.

\ivround{4}
\ivsays{The price is 100.5, and you must buy 2.}
\ivstop
\textbf{Model: 2, as you must.} You pay $201 + 2 = 203$.

\medskip
\textbf{Settlement.} The model paid $304.5 + 103 + 203 = 610.5$. The table compares a plain even split of 2, 2, 1, 1, and buying all 6 at once, and waiting to buy them all in round 4.
\begin{center}\small
\begin{tabular}{@{}lrr@{}}
 & Cost on this deal & Expected, with these forecasts \\ \hline
Model & $610.5$ & $607.5$ \\
Even split & $612.0$ & $609.0$ \\
All in round 1 & $618.0$ & $618.0$ \\
All in round 4 & $621.0$ & $615.0$
\end{tabular}
\end{center}
One deal proves little, though here the two agree. The noise lands on all of them, and only the lots still to buy are exposed to it, so the even splitter isn't punished for much. The gap comes from using the forecast.

\iqwrong{Treating the price level as a signal. A price that has risen isn't dear, and a price that has fallen isn't cheap: only the forecast says what comes next. The other error is to buy at once because you're anxious, which costs 18 of impact.}
\iqfollow{What if you dislike risk?}{You'd buy a bit sooner, since lots still to buy carry the price risk. With a large enough dislike, you'd buy much more in round 1. Say what you are trading off: certain impact against uncertain price.}
\iqfollow{Would a real interviewer give you these rules?}{They would rarely give a formula for impact, and in reality part of your impact stays in the price and shows others that you are buying. So you'd say what you assume, and how the answer changes if it stays.}
\iqwhy{Working an order is a daily decision on many desks. An interviewer watches whether you split the order and react to a forecast, and whether you avoid the urge to chase or wait for a dip.}
\end{solution}

\begin{solution}{iq:iv2:played-games:12}
Keep everything below this line covered, and slide the paper down one stop line at a time. The model comes from working backwards. You want the chance of finishing ahead, and nothing else. Say that first: an interviewer wants to know what you're trying to maximise before you bet.

\ivround{1}
\ivsays{You're 3 behind. A fair coin at even money. What's your stake?}
\ivstop
\textbf{Model: 4.} With a win, you're 1 ahead, and you stake nothing from then on. With a loss, you can still come back.
\begin{center}\small
\begin{tabular}{@{}l*{5}{r}@{}}
Stake & 0 & 1 & 2 & 3 & 4 \\ \hline
Chance of finishing ahead & 0.658 & 0.607 & 0.629 & 0.611 & 0.665
\end{tabular}
\end{center}
It's nearly a tie with staking nothing, and both beat small bets. Small bets are the worst here: they use up rounds without closing the gap. The coin loses, and you're 7 behind.

\ivround{2}
\ivsays{You're 7 behind. 55\% at even money. What's your stake?}
\ivstop
\textbf{Model: 4.} You must gain 8 points with stakes of at most 4, so you need two full wins, or to win all three rounds left. Stake 4 now, and a win leaves you 3 behind with the 60\% round to come: $0.55 \times 0.60 = 0.33$. No stake leaves rounds 3 and 4 at full stakes: $0.40 \times 0.60 = 0.24$. A stake of 1, 2 or 3 gives only 0.132, because you must then win all three. It wins, and you're 3 behind.

\ivround{3}
\ivsays{You're 3 behind. 40\% at even money, and one round to come at 60\%. What's your stake?}
\ivstop
\textbf{Model: 0.} If you stake nothing, one stake of 4 at 60\% next round wins it: a chance of 0.60. A stake of 4 here gives 0.40, since a win puts you ahead with nothing to do, and a loss ends it. Being behind doesn't mean taking any gamble: you take the best one that's left, and it's the next one. The coin loses, and you're still 3 behind.

\ivround{4}
\ivsays{You're 3 behind. 60\% at even money, and this is the last round. What's your stake?}
\ivstop
\textbf{Model: 4.} A stake of 3 ties, and a tie loses. Anything less than 4 is the same as nothing. It wins, and you finish 1 ahead.

\medskip
\textbf{Settlement.} The table compares the model with a player who maximises expected points, staking 4 at 55\% and 60\% and nothing else, and a player who stakes 1 every round.
\begin{center}\small
\begin{tabular}{@{}lrrr@{}}
 & Gap, this deal & Chance of finishing ahead & Expected final gap \\ \hline
Model & $+1$ & 0.665 & $-2.58$ \\
Most points & $+5$ & 0.330 & $-1.80$ \\
Stake 1 each round & $-3$ & 0.066 & $-2.90$
\end{tabular}
\end{center}
One deal proves little: both of the first two win here, and the one that plays for points has the bigger score. Over many deals, the model finishes ahead twice as often, and it pays for that with a worse expected gap: it gives up points to win the prize.

\iqwrong{Playing for points when the prize is for rank, or staking the maximum at every round because you're behind. In round 3, the stake that's right on ``I must gamble'' is wrong: better odds are still to come.}
\iqfollow{And if you were 1 ahead?}{Stake nothing, even at 60\%. You've already won, and a good bet can only lose you the prize. That's the other half of the lesson.}
\iqfollow{Would you trade like this on a desk?}{No. A bonus can feel like this prize, and that's exactly why a firm limits it. The firm wants expected profit for its risk, so say that you know the difference. And a real interviewer may be watching a game like this for whether you take silly risks: ask what the score means before you start.}
\iqwhy{Payoffs shape behaviour. An interviewer watches whether you ask what you're being scored on, and whether you can reason about risk when the payoff isn't linear.}
\end{solution}
